\section{Completion and change of ground field}\label{proj:sec:completion}

\begin{proof}[Proof of Theorem~\ref{proj:thm:main} over \(\C\)]
We prove existence of good log minimal models by dimension induction,
in the form used in Proposition~\ref{proj:prop:inductive-reduction}.
Dimension zero is immediate. Assume that good models exist in all
smaller dimensions. The signed-representative argument
(Theorem~\ref{proj:thm:signed}), the exclusions, and the jet estimates are
then available with exactly that lower-dimensional hypothesis.
Theorem~\ref{proj:thm:smooth-nonvanishing} establishes smooth
nonvanishing in the present dimension. Proposition~\ref{proj:prop:inductive-reduction}
therefore supplies the good-model statement in this dimension and
closes the induction.

Apply this to the rational lc pair \((X,B)\) of
Theorem~\ref{proj:thm:main}. On a common resolution of its dlt model and a
good log minimal model, the pullbacks of the adjoints agree because
the original adjoint is nef; this is the nef comparison in
Lemma~\ref{proj:mmp:comparison}. The pullback of \(K_X+B\)
is thus semiample. If \(p:W\to X\) is this resolution, normality gives
\(p_*\OO_W=\OO_X\). Choose a sufficiently divisible Cartier multiple
\(m(K_X+B)\) whose pullback is globally generated. The projection
formula identifies its sections with those of its pullback.
Surjectivity of evaluation upstairs implies surjectivity downstairs:
otherwise a base point downstairs would make every pulled-back
section vanish on its nonempty fiber. Hence \(K_X+B\) is semiample.
\end{proof}

\begin{proposition}[Change of algebraically closed field]\label{proj:prop:field-descent}
The conclusion of Theorem~\ref{proj:thm:main} over \(\C\) implies its
conclusion over every algebraically closed field of characteristic zero.
\end{proposition}

\begin{proof}
Let \((X,B)\) be defined over such a field \(k\). Choose a finitely
generated subfield \(k_1\subset k\) over which the projective variety,
the rational boundary, a Cartier multiple of its adjoint, and a log
resolution are all defined. After enlarging \(k_1\) if necessary,
these data base change to the given data. Let \(k_0\) be its algebraic
closure inside \(k\). It embeds into \(\C\).

The discrepancies on the chosen resolution are unchanged by
algebraically closed field extension. Since the resolution has simple
normal crossing boundary, it tests log canonicity over both \(k_0\)
and \(\C\).

Nefness is also invariant under such extensions. Indeed, a curve
over an extension is represented by a point of a relative Hilbert
scheme after its finite defining data have been spread over a
finite-type parameter scheme. The degree of a fixed line bundle is
constant in the resulting flat family. Specializing to a closed
point over the algebraically closed smaller field preserves a
negative degree, if one existed; some irreducible component of the
specialized curve would then have negative degree. This contradicts
nefness over that field. Conversely, curves over the smaller field
remain available after extension. Thus the \(k_0\)-model and its
complex base change are nef.

The complex case supplies an integer \(m>0\), enlarged to a multiple
of the Cartier index fixed over \(k_0\), for which the Cartier
line bundle \(M=\OO_{X_{k_0}}(m(K_{X_{k_0}}+B_{k_0}))\) becomes
globally generated over \(\C\). Proper flat base change for sections
identifies
\[
 H^0(X_{k_0},M)\otimes_{k_0}\C
   = H^0(X_{\C},M_{\C}).
\]
The cokernel of the evaluation map for \(M\) therefore becomes zero
after the faithfully flat extension \(k_0\subset\C\), and is already
zero over \(k_0\). Base change from \(k_0\) to \(k\) preserves this
surjectivity. The same integer \(m\) proves the required
semiampleness over \(k\).
\end{proof}

Together with Proposition~\ref{proj:prop:field-descent}, the complex proof
establishes Theorem~\ref{proj:thm:main} in its stated generality.
